% Presentacion - PFM02 / Sesion 03
% Multiplicacion y division con fracciones

\pdfminorversion=7
\documentclass[
	spanish,
	aspectratio=1610,
	xcolor={dvipsnames,table}
]{beamer}

\geometry{paperwidth=19.2cm,paperheight=12cm}

\usepackage[utf8]{inputenc}
\usepackage[T1]{fontenc}
\usepackage[spanish,es-nodecimaldot]{babel}
\usepackage{lmodern}
\usepackage{booktabs}
\usepackage{array}
\usepackage{graphicx}
\usepackage{ragged2e}
\usepackage{textcomp}
\usepackage{tikz}
\usetikzlibrary{arrows.meta,positioning,calc,decorations.pathreplacing,shapes.geometric,fit,backgrounds,overlay-beamer-styles}
\usepackage{hyperref}

% Datos del alumno y del programa
\newcommand{\studentname}{Mtro. Martin Jonathan de la Cruz Munoz}
\newcommand{\studentshort}{M. de la Cruz Munoz}
\newcommand{\studentid}{00841}
\newcommand{\studentemail}{martin.delacruz@campus.iiiepe.edu.mx}
\newcommand{\programname}{Maestria en Aprendizaje y Ensenanza de las Matematicas}
\newcommand{\programshort}{MAEM}
\newcommand{\coursecode}{PFM02}
\newcommand{\coursename}{Temas selectos de matematicas I}
\newcommand{\coursegroup}{Grupo A}
\newcommand{\courseperiod}{Marzo--julio 2026}
\newcommand{\institutionname}{IIIEPE}
\newcommand{\institutionlongname}{Instituto de Investigacion, Innovacion y Estudios de Posgrado para la Educacion del Estado de Nuevo Leon}
\newcommand{\institutionplace}{Monterrey, Nuevo Leon}
\newcommand{\iiiepelogo}{img/iiiepe/logo-iiiepe-blanco.png}
\newcommand{\iiiepemark}{img/iiiepe/marca-iiiepe.png}
\newcommand{\iiiepedots}{img/iiiepe/logo-puntos.png}

% Paleta institucional
\definecolor{iiiepeNavy}{HTML}{163B5C}
\definecolor{iiiepeBlue}{HTML}{1F6F9F}
\definecolor{iiiepeAqua}{HTML}{20B7A8}
\definecolor{iiiepeOrange}{HTML}{F28C28}
\definecolor{iiiepePaper}{HTML}{F6F8F7}
\definecolor{iiiepeInk}{HTML}{1E2A32}
\definecolor{iiiepeMuted}{HTML}{667085}

\mode<presentation>{
	\usetheme{default}
	\usefonttheme{professionalfonts}
	\setbeamertemplate{navigation symbols}{}
	\setbeamertemplate{blocks}[rounded][shadow=false]
}

\hypersetup{
	colorlinks=true,
	urlcolor=iiiepeBlue,
	linkcolor=iiiepeNavy
}

\setbeamercolor{background canvas}{bg=white}
\setbeamercolor{normal text}{fg=iiiepeInk,bg=white}
\setbeamercolor{structure}{fg=iiiepeBlue}
\setbeamercolor{frametitle}{fg=white,bg=iiiepeNavy}
\setbeamercolor{title}{fg=white,bg=iiiepeNavy}
\setbeamercolor{block title}{fg=white,bg=iiiepeBlue}
\setbeamercolor{block body}{fg=iiiepeInk,bg=iiiepePaper}
\setbeamercolor{alerted text}{fg=iiiepeOrange}
\setbeamerfont{title}{size=\LARGE,series=\bfseries}
\setbeamerfont{frametitle}{size=\large,series=\bfseries}
\setbeamerfont{block title}{series=\bfseries}
\setbeamertemplate{itemize item}{\textcolor{iiiepeAqua}{\large$\blacktriangleright$}}
\setbeamertemplate{itemize subitem}{\textcolor{iiiepeOrange}{\scriptsize$\blacksquare$}}
\setbeamertemplate{enumerate item}{\textcolor{iiiepeBlue}{\insertenumlabel.}}

\newcommand{\localLogo}[1]{%
	\IfFileExists{#1}{\includegraphics[height=0.58cm]{#1}}{\textbf{\institutionname}}%
}

\setbeamertemplate{frametitle}{
	\nointerlineskip
	\begin{beamercolorbox}[wd=\textwidth,ht=1.05cm,dp=0.22cm,leftskip=0.35cm,rightskip=0.25cm]{frametitle}
		\begin{minipage}[c]{0.72\textwidth}
			\insertframetitle
		\end{minipage}
		\hfill
		\raisebox{-0.1cm}{\localLogo{\iiiepelogo}}
	\end{beamercolorbox}
	\vspace{-0.04cm}
	{\color{iiiepeOrange}\rule{\textwidth}{1.4pt}}
}

\setbeamertemplate{footline}{
	\leavevmode
	\hbox{%
		\begin{beamercolorbox}[wd=.34\paperwidth,ht=2.7ex,dp=1ex,leftskip=1em]{author in head/foot}
			\color{white}\studentshort
		\end{beamercolorbox}
		\begin{beamercolorbox}[wd=.38\paperwidth,ht=2.7ex,dp=1ex,center]{title in head/foot}
			\color{white}\programshort\ -- \coursecode
		\end{beamercolorbox}
		\begin{beamercolorbox}[wd=.20\paperwidth,ht=2.7ex,dp=1ex,rightskip=1em plus 1fil]{date in head/foot}
			\color{white}Matricula \studentid\hfill\insertframenumber/\inserttotalframenumber
		\end{beamercolorbox}
	}
}
\setbeamercolor{author in head/foot}{bg=iiiepeNavy}
\setbeamercolor{title in head/foot}{bg=iiiepeBlue}
\setbeamercolor{date in head/foot}{bg=iiiepeNavy}

\newcommand{\accentbar}{%
	\begin{tikzpicture}[remember picture,overlay]
		\path[use as bounding box] (0,0) rectangle (0,0);
		\fill[iiiepeAqua] (current page.south west) rectangle ([xshift=0.42\paperwidth,yshift=0.08cm]current page.south west);
		\fill[iiiepeOrange] ([xshift=0.42\paperwidth]current page.south west) rectangle ([xshift=\paperwidth,yshift=0.08cm]current page.south west);
	\end{tikzpicture}
}

\newcommand{\DividirEntre}[2]{\ensuremath{#1\,\text{ entre }\,#2}}
\newcommand{\signodiv}{\mathbin{\text{\normalfont\textdiv}}}

\tikzset{
	cajaProc/.style={
		draw=iiiepeBlue,
		rounded corners=2mm,
		align=center,
		inner sep=4pt,
		minimum height=8mm,
		font=\small
	},
	cajaDato/.style={
		cajaProc,
		fill=iiiepeBlue!8
	},
	cajaOp/.style={
		cajaProc,
		fill=iiiepeOrange!12
	},
	cajaRes/.style={
		cajaProc,
		fill=iiiepeAqua!12
	},
	flechaProc/.style={
		-{Latex[length=2.4mm]},
		thick,
		draw=iiiepeBlue
	},
	notaProc/.style={
		draw=iiiepeOrange,
		rounded corners=2mm,
		fill=iiiepeOrange!12,
		align=center,
		inner sep=4pt,
		font=\scriptsize
	},
	decisionProc/.style={
		diamond,
		aspect=2.1,
		draw=iiiepeBlue,
		fill=iiiepeAqua!10,
		align=center,
		inner sep=2pt,
		font=\scriptsize
	}
}

\newcommand{\RutaOperacion}[4]{%
\begin{tikzpicture}[node distance=6mm]
	\node[cajaDato,visible on=<1->] (a) {#1};
	\node[cajaOp, right=of a,visible on=<2->] (b) {#2};
	\node[cajaOp, right=of b,visible on=<3->] (c) {#3};
	\node[cajaRes, right=of c,visible on=<4->] (d) {#4};
	\draw[flechaProc,visible on=<1->] (a) -- (b);
	\draw[flechaProc,visible on=<2->] (b) -- (c);
	\draw[flechaProc,visible on=<3->] (c) -- (d);
\end{tikzpicture}%
}

\newcommand{\RutaOperacionCinco}[5]{%
\begin{tikzpicture}[node distance=5.5mm]
	\node[cajaDato,visible on=<1->] (a) {#1};
	\node[cajaOp, right=of a,visible on=<2->] (b) {#2};
	\node[cajaOp, right=of b,visible on=<3->] (c) {#3};
	\node[cajaOp, right=of c,visible on=<4->] (d) {#4};
	\node[cajaRes, right=of d,visible on=<5->] (e) {#5};
	\draw[flechaProc,visible on=<1->] (a) -- (b);
	\draw[flechaProc,visible on=<2->] (b) -- (c);
	\draw[flechaProc,visible on=<3->] (c) -- (d);
	\draw[flechaProc,visible on=<4->] (d) -- (e);
\end{tikzpicture}%
}

\newcommand{\GridProducto}[5]{%
\begin{tikzpicture}[scale=#5]
	\pgfmathsetmacro{\yini}{#4-#3}
	\pgfmathsetmacro{\xmitad}{#1/2}
	\pgfmathsetmacro{\ymitad}{#4-#3/2}
	\pgfmathsetmacro{\yarriba}{#4+0.22}
	\pgfmathsetmacro{\xderecha}{#2+0.18}
	\draw[step=1,gray!55,very thin,visible on=<1->] (0,0) grid (#2,#4);
	\draw[thick,visible on=<1->] (0,0) rectangle (#2,#4);
	\fill[iiiepeBlue!18,visible on=<2->] (0,0) rectangle (#1,#4);
	\fill[iiiepeOrange!18,visible on=<3->] (0,\yini) rectangle (#2,#4);
	\fill[iiiepeAqua!55,visible on=<4->] (0,\yini) rectangle (#1,#4);
	\draw[decorate,decoration={brace,mirror,raise=3pt},visible on=<5->] (0,-0.08) -- (#2,-0.08)
	node[midway,below=6pt] {\scriptsize #2 partes};
	\draw[decorate,decoration={brace,raise=3pt},visible on=<5->] (-0.08,0) -- (-0.08,#4)
	node[midway,left=6pt] {\scriptsize #4 partes};
	\node[above,visible on=<5->] at (\xmitad,\yarriba) {\scriptsize #1 tomadas};
	\node[right,visible on=<5->] at (\xderecha,\ymitad) {\scriptsize #3 tomadas};
\end{tikzpicture}%
}

\newcommand{\sectionframe}[1]{%
	\begin{frame}[plain]
		\begin{tikzpicture}[remember picture,overlay]
			\fill[iiiepeNavy] (current page.south west) rectangle (current page.north east);
			\IfFileExists{\iiiepedots}{%
				\node[opacity=0.13,anchor=east] at ([xshift=0.8cm]current page.east) {\includegraphics[height=9.4cm]{\iiiepedots}};
			}{}
			\fill[iiiepeAqua] ([xshift=0.85cm,yshift=2.1cm]current page.south west) rectangle ([xshift=1.05cm,yshift=6.3cm]current page.south west);
			\fill[iiiepeOrange] ([xshift=1.15cm,yshift=2.1cm]current page.south west) rectangle ([xshift=1.35cm,yshift=5.35cm]current page.south west);
		\end{tikzpicture}
		\vspace{1.8cm}
		\hspace{1.25cm}
		\begin{minipage}{0.72\paperwidth}
			{\color{white}\Huge\bfseries #1}\\[0.25cm]
			{\color{white!75}Sesion 03 -- Multiplicacion y division con fracciones}
		\end{minipage}
	\end{frame}
}

% Separadores automaticos desactivados para evitar diapositivas repetitivas.
% \AtBeginSection[]{\sectionframe{\insertsectionhead}}

\title[\coursecode]{Multiplicacion y division con fracciones}
\subtitle{Sesion 03 -- \coursename}
\author[\studentshort]{\texorpdfstring{\studentname\\\footnotesize Matricula: \studentid\\\href{mailto:\studentemail}{\studentemail}}{\studentname, Matricula: \studentid}}
\institute[\institutionname]{\institutionlongname\\\coursecode\ -- \coursegroup}
\date[\courseperiod]{\texorpdfstring{\institutionplace\\\courseperiod}{\institutionplace, \courseperiod}}

\begin{document}

\begin{frame}[plain]
	\begin{tikzpicture}[remember picture,overlay]
		\fill[iiiepeNavy] (current page.south west) rectangle ([xshift=0.56\paperwidth]current page.north west);
		\fill[iiiepeBlue] ([xshift=0.56\paperwidth]current page.south west) rectangle (current page.north east);
		\IfFileExists{\iiiepedots}{%
			\node[opacity=0.16,anchor=south east] at ([xshift=0.8cm,yshift=-0.3cm]current page.south east) {\includegraphics[height=8.6cm]{\iiiepedots}};
		}{}
		\node[anchor=north west] at ([xshift=0.55cm,yshift=-0.45cm]current page.north west) {\localLogo{\iiiepelogo}};
		\fill[iiiepeOrange] ([xshift=0.58\paperwidth,yshift=1.20cm]current page.south west) rectangle ([xshift=0.92\paperwidth,yshift=1.34cm]current page.south west);
		\fill[iiiepeAqua] ([xshift=0.58\paperwidth,yshift=0.90cm]current page.south west) rectangle ([xshift=0.82\paperwidth,yshift=1.04cm]current page.south west);
	\end{tikzpicture}

	\vspace{1.62cm}
	\begin{minipage}{0.52\paperwidth}
		{\color{white}\Huge\bfseries Multiplicacion y division\\ con fracciones}\\[0.25cm]
		{\color{white!86}\large PFM02 -- Sesion 03}\\[0.35cm]
		{\color{iiiepeOrange}\rule{0.82\linewidth}{1.3pt}}\\[0.45cm]
		{\color{white}\studentname}\\
		{\color{white!80}\footnotesize Matricula: \studentid}
	\end{minipage}
	\hfill
	\begin{minipage}{0.35\paperwidth}
		\raggedleft
		{\color{white}\Large\bfseries \institutionname}\\[0.18cm]
		{\color{white!82}\footnotesize \programname}\\[0.5cm]
		{\color{white}\small \coursecode\ -- \coursename}\\
		{\color{white!78}\small \courseperiod}\\[0.5cm]
		{\color{white!74}\footnotesize \institutionplace}
	\end{minipage}
\end{frame}

\begin{frame}{Contenido}
	\tableofcontents
\end{frame}

\section{Actividad}

\begin{frame}{Ficha de la sesion}
	\begin{columns}[T]
		\begin{column}{0.48\textwidth}
			\begin{block}{Contenido}
				Extension del significado de las operaciones y sus relaciones inversas.
			\end{block}
			\begin{block}{PDA}
				Reconoce el significado de las cuatro operaciones basicas y sus relaciones inversas al resolver problemas con numeros con signo.
			\end{block}
		\end{column}
		\begin{column}{0.48\textwidth}
			\begin{block}{Aprendizaje de la secuencia}
				Resolver situaciones problematicas que involucran multiplicar o dividir fracciones.
			\end{block}
			\begin{block}{Producto}
				Procedimientos visibles, ejercicios resueltos, situaciones en contexto, evaluacion y retroalimentacion.
			\end{block}
		\end{column}
	\end{columns}
	\accentbar
\end{frame}

\begin{frame}{Criterios incorporados}
	\small
	\begin{columns}[T]
		\begin{column}{0.49\textwidth}
			\begin{block}{Estructura de secuencia}
				La actividad se organiza en cinco fases: antecedente, explicacion, practica, evaluacion y retroalimentacion.
			\end{block}
			\begin{block}{Evidencia esperada}
				No basta la respuesta: debe observarse estrategia, procedimiento, simplificacion e interpretacion.
			\end{block}
		\end{column}
		\begin{column}{0.49\textwidth}
			\begin{block}{Criterios de evaluacion}
				Actividades diarias, exposicion de secuencia, escrito reflexivo y exposicion final requieren claridad matematica y reflexion.
			\end{block}
			\begin{block}{Decision didactica}
				Cada ejemplo muestra que se hace, por que funciona y como se comprueba.
			\end{block}
		\end{column}
	\end{columns}
	\accentbar
\end{frame}

\begin{frame}{Ruta pedagogica de la actividad}
	\begin{columns}[T]
		\begin{column}{0.20\textwidth}
			\begin{block}{1. Recuperar}
				Signos, simplificacion, enteros y mixtas.
			\end{block}
		\end{column}
		\begin{column}{0.20\textwidth}
			\begin{block}{2. Interpretar}
				Producto: parte de parte. Division: cuantos caben.
			\end{block}
		\end{column}
		\begin{column}{0.20\textwidth}
			\begin{block}{3. Transformar}
				Usar impropias y reciproco cuando hay division.
			\end{block}
		\end{column}
		\begin{column}{0.20\textwidth}
			\begin{block}{4. Operar}
				Multiplicar numeradores y denominadores.
			\end{block}
		\end{column}
		\begin{column}{0.20\textwidth}
			\begin{block}{5. Verificar}
				Simplificar, revisar signo y unidad.
			\end{block}
		\end{column}
	\end{columns}
	\vspace{0.25cm}
	\begin{block}{Criterio docente}
		El procedimiento es suficiente cuando se entiende que multiplicar por una fraccion puede reducir y dividir entre una fraccion pregunta cuantas veces cabe esa parte.
	\end{block}
	\accentbar
\end{frame}

\begin{frame}{Mapa general de multiplicacion y division}
	\centering
	\resizebox{0.96\textwidth}{!}{%
	\begin{tikzpicture}[node distance=0.85cm and 1.0cm]
		\node[cajaDato] (inicio) {Fracciones};
		\node[decisionProc, right=of inicio] (op) {operacion};
		\node[cajaOp, above right=0.45cm and 1.0cm of op] (mult) {multiplicacion:\\parte de parte};
		\node[cajaOp, below right=0.45cm and 1.0cm of op] (divi) {division:\\usar reciproco};
		\node[cajaOp, right=2.7cm of op] (convertir) {convertir\\mixtas o enteros};
		\node[cajaOp, right=of convertir] (operar) {multiplicar\\factores};
		\node[cajaOp, right=of operar] (simplificar) {simplificar};
		\node[cajaRes, right=of simplificar] (verificar) {interpretar\\unidad};
		\draw[flechaProc] (inicio) -- (op);
		\draw[flechaProc] (op) -- (mult);
		\draw[flechaProc] (op) -- (divi);
		\draw[flechaProc] (mult) -- (convertir);
		\draw[flechaProc] (divi) -- (convertir);
		\draw[flechaProc] (convertir) -- (operar);
		\draw[flechaProc] (operar) -- (simplificar);
		\draw[flechaProc] (simplificar) -- (verificar);
	\end{tikzpicture}
	}
	\vspace{0.25cm}
	\begin{block}{Lectura del mapa}
		La division se transforma en multiplicacion por el reciproco; despues ambos procedimientos se revisan con simplificacion, signo y sentido contextual.
	\end{block}
	\accentbar
\end{frame}

\begin{frame}{Idea clave: parte de parte y cuantos caben}
	\begin{columns}[T]
		\begin{column}{0.48\textwidth}
			\begin{block}{Multiplicar fracciones}
				Una fraccion de otra fraccion se representa como superposicion de partes.
				\[
					\frac{a}{b}\times\frac{c}{d}=\frac{ac}{bd}
				\]
			\end{block}
			\centering
			\GridProducto{3}{4}{5}{9}{0.30}
		\end{column}
		\begin{column}{0.48\textwidth}
			\begin{block}{Dividir fracciones}
				Dividir entre una fraccion equivale a contar grupos de ese tamano.
				\[
					\frac{a}{b}\signodiv\frac{c}{d}=\frac{a}{b}\times\frac{d}{c}
				\]
			\end{block}
			\centering
			\resizebox{0.98\linewidth}{!}{%
			\RutaOperacion
			{$\frac{7}{8}\signodiv\frac{1}{6}$}
			{$\frac{7}{8}\times\frac{6}{1}$}
			{$\frac{42}{8}$}
			{$5\,\frac14$}}
		\end{column}
	\end{columns}
	\accentbar
\end{frame}

\section{Procedimientos}

\begin{frame}{Procedimiento 1: multiplicar fracciones}
	\small
	\begin{columns}[T]
		\begin{column}{0.48\textwidth}
			\begin{block}{Regla operativa}
				\[
					\frac{3}{4}\times\frac{5}{9}
					=\frac{3\cdot5}{4\cdot9}
					=\frac{15}{36}
					=\frac{5}{12}
				\]
			\end{block}
			\begin{block}{Sentido}
				Se toma una parte de otra parte. La zona doblemente sombreada representa el numerador del producto.
			\end{block}
		\end{column}
		\begin{column}{0.48\textwidth}
			\centering
			\GridProducto{3}{4}{5}{9}{0.39}
			\vspace{0.15cm}
			\[
				\text{interseccion}=\frac{3\cdot5}{4\cdot9}
			\]
		\end{column}
	\end{columns}
	\begin{block}{Verificacion}
		Despues de multiplicar se revisa si numerador y denominador comparten divisor comun.
	\end{block}
	\accentbar
\end{frame}

\begin{frame}{Procedimiento 2: enteros y fracciones mixtas}
	\small
	\begin{columns}[T]
		\begin{column}{0.48\textwidth}
			\begin{block}{Convertir antes de operar}
				\[
					3\,\frac25=\frac{3(5)+2}{5}=\frac{17}{5},
					\qquad
					1\,\frac17=\frac{8}{7}
				\]
			\end{block}
			\begin{block}{Entero como fraccion}
				\[
					7\times\frac{3}{11}
					=\frac{7}{1}\times\frac{3}{11}
					=\frac{21}{11}
				\]
			\end{block}
		\end{column}
		\begin{column}{0.48\textwidth}
			\begin{block}{Ejemplo completo}
				\[
					3\,\frac25\times1\,\frac17
					=\frac{17}{5}\times\frac{8}{7}
					=\frac{136}{35}
					=3\,\frac{31}{35}
				\]
			\end{block}
			\centering
			\resizebox{0.98\linewidth}{!}{%
			\RutaOperacionCinco
			{$3\,\frac25$}
			{$\frac{17}{5}$}
			{$\frac{17}{5}\times\frac87$}
			{$\frac{136}{35}$}
			{$3\,\frac{31}{35}$}}
		\end{column}
	\end{columns}
	\accentbar
\end{frame}

\begin{frame}{Procedimiento 3: dividir por el reciproco}
	\small
	\begin{columns}[T]
		\begin{column}{0.48\textwidth}
			\begin{block}{Regla}
				\[
					\frac{a}{b}\signodiv\frac{c}{d}
					=
					\frac{a}{b}\times\frac{d}{c}
				\]
			\end{block}
			\begin{block}{Ejemplo}
				\[
					\frac{7}{8}\signodiv\frac{1}{6}
					=
					\frac{7}{8}\times\frac{6}{1}
					=
					\frac{42}{8}
					=
					5\,\frac14
				\]
			\end{block}
		\end{column}
		\begin{column}{0.48\textwidth}
			\begin{alertblock}{Cuidado}
				Solo se invierte la segunda fraccion: la que indica el tamano del grupo o divisor.
			\end{alertblock}
			\centering
			\resizebox{0.95\linewidth}{!}{%
			\RutaOperacion
			{$\frac{7}{8}\signodiv\frac{1}{6}$}
			{$\frac{7}{8}\times\frac{6}{1}$}
			{$\frac{42}{8}$}
			{$5\,\frac14$}}
		\end{column}
	\end{columns}
	\begin{block}{Sentido del atajo}
		Dividir por $\frac{1}{6}$ significa preguntar cuantas sextas partes caben en $\frac{7}{8}$.
	\end{block}
	\accentbar
\end{frame}

\begin{frame}{Procedimiento 4: signos y simplificacion}
	\small
	\begin{columns}[T]
		\begin{column}{0.48\textwidth}
			\begin{block}{Signos}
				\[
					(+)(+)=+,
					\qquad (+)(-)=-,
					\qquad (-)(-) = +
				\]
			\end{block}
			\begin{block}{Ejemplo con signo}
				\[
					-\frac{4}{7}\times\frac{3}{6}
					=
					-\frac{12}{42}
					=
					-\frac{2}{7}
				\]
			\end{block}
		\end{column}
		\begin{column}{0.48\textwidth}
			\begin{block}{Simplificar}
				\[
					\frac{5}{6}\signodiv\frac{2}{3}
					=
					\frac{5}{6}\times\frac{3}{2}
					=
					\frac{15}{12}
					=
					\frac54
				\]
			\end{block}
			\begin{alertblock}{Revision obligatoria}
				El resultado debe quedar simplificado y, si conviene, expresado como mixta para interpretarlo.
			\end{alertblock}
		\end{column}
	\end{columns}
	\centering
	\resizebox{0.66\textwidth}{!}{%
	\RutaOperacion
	{$\frac{15}{12}$}
	{MCD $3$}
	{$\frac{15/3}{12/3}$}
	{$\frac54=1\,\frac14$}}
	\accentbar
\end{frame}

\section{Ejercicios resueltos}

\begin{frame}{Ejercicios: multiplicaciones basicas}
	\small
	La tabla presenta varios incisos resueltos; el esquema posterior muestra la ruta comun que se aplica al bloque.
	\vspace{0.1cm}
	\begin{block}{Multiplicar factores y simplificar}
		\centering
		\begin{tabular}{p{0.21\textwidth}p{0.43\textwidth}p{0.19\textwidth}}
			\toprule
			\textbf{Inciso} & \textbf{Operacion} & \textbf{Resultado}\\
			\midrule
			a & $\frac34\times\frac59=\frac{15}{36}=\frac{5}{12}$ & $\frac{5}{12}$\\
			b & $\frac49\cdot\frac{12}{7}=\frac{48}{63}=\frac{16}{21}$ & $\frac{16}{21}$\\
			c & $\frac{23}{25}\cdot\frac65=\frac{138}{125}$ & $1\,\frac{13}{125}$\\
			\bottomrule
		\end{tabular}
	\end{block}
	\centering
	\resizebox{0.82\linewidth}{!}{%
	\RutaOperacion
	{$\frac34\times\frac59$}
	{$\frac{3\cdot5}{4\cdot9}$}
	{$\frac{15}{36}$}
	{$\frac{5}{12}$}}
	\begin{block}{Criterio}
		El producto no siempre es mayor que los factores: si se multiplica por una fraccion menor que $1$, la cantidad puede disminuir.
	\end{block}
	\accentbar
\end{frame}

\begin{frame}{Ejercicios: mixtas y factores}
	\small
	La tabla presenta varios incisos resueltos; el esquema posterior muestra la ruta comun que se aplica al bloque.
	\begin{columns}[T]
		\begin{column}{0.49\textwidth}
			\begin{block}{Resultados}
				\scriptsize
				\setlength{\tabcolsep}{2pt}
				\begin{tabular}{@{}p{0.16\linewidth}p{0.56\linewidth}p{0.20\linewidth}@{}}
					\toprule
					\textbf{Inciso} & \textbf{Operacion} & \textbf{Resultado}\\
					\midrule
					d & $7\times\frac{3}{11}=\frac{21}{11}$ & $1\,\frac{10}{11}$\\
					e & $3\,\frac25\times1\,\frac17=\frac{17}{5}\times\frac87$ & $3\,\frac{31}{35}$\\
					f & $\frac34\times\frac25\times\frac69=\frac{36}{180}$ & $\frac15$\\
					\bottomrule
				\end{tabular}
			\end{block}
		\end{column}
		\begin{column}{0.41\textwidth}
			\centering
			\resizebox{0.92\linewidth}{!}{%
			\RutaOperacionCinco
			{$3\,\frac25$}
			{$\frac{17}{5}$}
			{$\frac{17}{5}\times\frac87$}
			{$\frac{136}{35}$}
			{$3\,\frac{31}{35}$}}
			\vspace{0.15cm}
			\begin{alertblock}{Cuidado}
				Convierte la mixta completa antes de operar.
			\end{alertblock}
		\end{column}
	\end{columns}
	\accentbar
\end{frame}

\begin{frame}{Ejercicios: divisiones por reciproco I}
	\small
	La tabla presenta varios incisos resueltos; el esquema posterior muestra la ruta comun que se aplica al bloque.
	\begin{block}{Cambiar division por multiplicacion}
		\centering
		\begin{tabular}{p{0.18\textwidth}p{0.48\textwidth}p{0.18\textwidth}}
			\toprule
			\textbf{Inciso} & \textbf{Operacion} & \textbf{Resultado}\\
			\midrule
			a & $\frac78\signodiv\frac16=\frac78\times\frac61=\frac{42}{8}$ & $5\,\frac14$\\
			b & $\frac57\signodiv\frac13=\frac57\times\frac31=\frac{15}{7}$ & $2\,\frac17$\\
			c & $\frac27\signodiv\frac67=\frac27\times\frac76=\frac{14}{42}$ & $\frac13$\\
			\bottomrule
		\end{tabular}
	\end{block}
	\centering
	\resizebox{0.72\textwidth}{!}{%
	\RutaOperacion
	{$\frac78\signodiv\frac16$}
	{$\frac78\times\frac61$}
	{$\frac{42}{8}$}
	{$5\,\frac14$}}
	\vspace{0.2cm}
	\begin{block}{Revision}
		En division, invertir la primera fraccion cambia el problema; solo se invierte el divisor.
	\end{block}
	\accentbar
\end{frame}

\begin{frame}{Ejercicios: divisiones con mixtas y simplificacion}
	\small
	La tabla presenta varios incisos resueltos; el esquema posterior muestra la ruta comun que se aplica al bloque.
	\begin{block}{Convertir, usar reciproco, simplificar}
		\centering
		\scriptsize
		\begin{tabular}{p{0.15\textwidth}p{0.55\textwidth}p{0.16\textwidth}}
			\toprule
			\textbf{Inciso} & \textbf{Operacion} & \textbf{Resultado}\\
			\midrule
			d & $1\,\frac27\signodiv2\,\frac58=\frac97\signodiv\frac{21}{8}=\frac97\times\frac8{21}$ & $\frac{24}{49}$\\
			e & $\frac56\signodiv\frac23=\frac56\times\frac32=\frac{15}{12}$ & $1\,\frac14$\\
			f & $\frac34\signodiv\frac25=\frac34\times\frac52$ & $1\,\frac78$\\
			g & $\frac89\signodiv\frac32=\frac89\times\frac23$ & $\frac{16}{27}$\\
			\bottomrule
		\end{tabular}
	\end{block}
	\centering
	\resizebox{0.74\textwidth}{!}{%
	\RutaOperacionCinco
	{$1\,\frac27\signodiv2\,\frac58$}
	{$\frac97\signodiv\frac{21}{8}$}
	{$\frac97\times\frac8{21}$}
	{$\frac{72}{147}$}
	{$\frac{24}{49}$}}
	\accentbar
\end{frame}

\section{Situaciones}

\begin{frame}{Situaciones: multiplicar una parte}
	\small
	\begin{columns}[T]
		\begin{column}{0.48\textwidth}
			\begin{block}{02b Tela}
				Usar la mitad de $\frac34$ m$^2$:
				\[
					\frac12\times\frac34=\frac38.
				\]
			\end{block}
			\begin{block}{03 Deuda}
				\[
					\frac34(132)=99,
					\qquad
					\frac23(99)=66,
					\qquad
					99-66=33.
				\]
			\end{block}
		\end{column}
		\begin{column}{0.48\textwidth}
			\begin{block}{10 Queso}
				\[
					\frac34\times4\,\frac23
					=
					\frac34\times\frac{14}{3}
					=
					\frac{42}{12}=\frac72.
				\]
			\end{block}
			\centering
			\GridProducto{3}{4}{1}{2}{0.62}
		\end{column}
	\end{columns}
	\accentbar
\end{frame}

\begin{frame}{Situaciones: dividir entre una medida}
	\small
	\begin{columns}[T]
		\begin{column}{0.48\textwidth}
			\begin{block}{04 Varillas}
				\[
					4\,\frac34\signodiv\frac14
					=
					\frac{19}{4}\times\frac41=19.
				\]
			\end{block}
			\begin{block}{08 Botellas}
				\[
					32\,\frac15\signodiv1\,\frac25
					=
					\frac{161}{5}\signodiv\frac75=23.
				\]
			\end{block}
		\end{column}
		\begin{column}{0.48\textwidth}
			\begin{block}{09 Vasos}
				\[
					11\,\frac14\signodiv\frac34
					=
					\frac{45}{4}\times\frac43=15.
				\]
			\end{block}
			\begin{block}{11 Bolsas}
				\[
					1\,\frac12\signodiv\frac{1}{10}
					=
					\frac32\times10=15.
				\]
			\end{block}
		\end{column}
	\end{columns}
	\accentbar
\end{frame}

\begin{frame}{Situaciones: area y medida faltante}
	\small
	\begin{columns}[T]
		\begin{column}{0.46\textwidth}
			\begin{block}{06 Ancho de mesa}
				\[
					ancho=\frac{area}{largo}
					=
					\frac{21}{4}\signodiv\frac72
					=
					\frac32\text{ m}.
				\]
			\end{block}
			\begin{block}{07 Piso de bano}
				\[
					A=\frac{13}{7}\times\frac45
					=
					\frac{52}{35}
					=
					1\,\frac{17}{35}.
				\]
			\end{block}
		\end{column}
		\begin{column}{0.50\textwidth}
			\centering
			\begin{tikzpicture}[scale=0.78]
				\draw[thick, fill=iiiepeAqua!12] (0,0) rectangle (4.6,2.0);
				\draw[Latex-Latex] (0,-0.35) -- (4.6,-0.35) node[midway,below] {$3\,\frac12$ m};
				\draw[Latex-Latex] (-0.35,0) -- (-0.35,2.0) node[midway,left] {$?$};
				\node at (2.3,1.0) {$A=\frac{21}{4}\text{ m}^2$};
				\node[cajaOp, text width=3.8cm] at (7.4,1.0) {
					$\frac{21}{4}\signodiv\frac72$\\[1mm]
					$=\frac{21}{4}\times\frac27$\\[1mm]
					$=\frac32$ m
				};
			\end{tikzpicture}
		\end{column}
	\end{columns}
	\accentbar
\end{frame}

\begin{frame}{Situaciones: procesos encadenados}
	\small
	\begin{columns}[T]
		\begin{column}{0.48\textwidth}
			\begin{block}{01 Desperdicio alimentario}
				\[
					8\times\frac13=\frac83=2\,\frac23,
					\qquad
					\frac12\times\frac13=\frac16.
				\]
			\end{block}
			\begin{block}{02a Jugos}
				\[
					3\times12\times\frac34=36\times\frac34=27\text{ L}.
				\]
			\end{block}
		\end{column}
		\begin{column}{0.48\textwidth}
			\begin{block}{05 Finca}
				\[
					\frac25(20)=8,
					\qquad
					20-8=12,
					\qquad
					\frac34(12)=9.
				\]
				Quedan $3$ hectareas.
			\end{block}
			\centering
			\resizebox{0.95\linewidth}{!}{%
			\RutaOperacionCinco
			{$20$ ha}
			{vender\\$8$}
			{restan\\$12$}
			{alquilar\\$9$}
			{quedan\\$3$}}
		\end{column}
	\end{columns}
	\accentbar
\end{frame}

\section{Evaluacion}

\begin{frame}{Autoevaluacion y evidencia}
	\scriptsize
	\begin{tabular}{p{0.58\textwidth}ccc}
		\toprule
		\textbf{Indicador} & \textbf{3} & \textbf{2} & \textbf{1}\\
		\midrule
		Realizo multiplicaciones de fracciones usando el algoritmo aprendido. & $\square$ & $\square$ & $\square$\\[0.16cm]
		Realizo divisiones de fracciones usando el reciproco de forma correcta. & $\square$ & $\square$ & $\square$\\[0.16cm]
		Identifico la operacion necesaria para resolver un problema. & $\square$ & $\square$ & $\square$\\[0.16cm]
		Convierto mixtas, simplifico e interpreto la respuesta contextual. & $\square$ & $\square$ & $\square$\\
		\bottomrule
	\end{tabular}
	\vspace{0.25cm}
	\begin{block}{Cierre metacognitivo}
		Que error debo vigilar: invertir la fraccion equivocada, no convertir mixtas, olvidar unidades o creer que todo producto aumenta.
	\end{block}
	\accentbar
\end{frame}

\begin{frame}{Retroalimentacion}
	\small
	\begin{columns}[T]
		\begin{column}{0.49\textwidth}
			\begin{block}{Errores frecuentes}
				\begin{itemize}
					\item Multiplicar sin simplificar el resultado.
					\item Dividir sin usar el reciproco.
					\item Invertir la primera fraccion en vez del divisor.
					\item Tratar la mixta como entero separado.
				\end{itemize}
			\end{block}
		\end{column}
		\begin{column}{0.49\textwidth}
			\begin{block}{Preguntas de revision}
				\begin{itemize}
					\item Estoy tomando una parte de otra parte?
					\item Estoy preguntando cuantos grupos caben?
					\item Converti enteros y mixtas a fracciones?
					\item La unidad final responde al contexto?
				\end{itemize}
			\end{block}
		\end{column}
	\end{columns}
	\accentbar
\end{frame}

\section{Cierre}

\begin{frame}{Conclusiones}
	\begin{enumerate}
		\item Multiplicar fracciones significa tomar una parte de otra parte.
		\item Dividir entre una fraccion significa contar cuantos grupos de ese tamano caben en una cantidad.
		\item El reciproco no es una regla aislada: expresa la relacion inversa entre multiplicacion y division.
		\item Las respuestas deben simplificarse, conservar signo correcto e interpretarse con unidad.
		\item La actividad sigue las cinco fases institucionales: antecedente, explicacion, practica, evaluacion y retroalimentacion.
	\end{enumerate}
	\accentbar
\end{frame}

\begin{frame}[plain]
	\begin{tikzpicture}[remember picture,overlay]
		\fill[iiiepeNavy] (current page.south west) rectangle (current page.north east);
		\IfFileExists{\iiiepedots}{%
			\node[opacity=0.13,anchor=south east] at ([xshift=0.8cm,yshift=-0.2cm]current page.south east) {\includegraphics[height=9cm]{\iiiepedots}};
		}{}
		\node[anchor=north west] at ([xshift=0.55cm,yshift=-0.45cm]current page.north west) {\localLogo{\iiiepelogo}};
	\end{tikzpicture}
	\vspace{2.2cm}
	\begin{center}
		{\color{white}\Huge\bfseries Gracias}\\[0.55cm]
		{\color{white}\Large \studentname}\\[0.18cm]
		{\color{white!78}\footnotesize Matricula \studentid\ -- \studentemail}\\[0.45cm]
		{\color{iiiepeOrange}\rule{0.36\paperwidth}{1.4pt}}\\[0.35cm]
		{\color{white!72}\footnotesize \institutionname\ -- \coursecode\ -- \coursename}
	\end{center}
\end{frame}

\end{document}

